% Open ball, closed ball, and a neighborhood of a point.
%
%   figures/tikz/ra-balls.tex
%       -> media/figures/ra-balls.{png,svg}
%
% Lecture 13 (Real Analysis Review), Sets and sequences.
%
% GEOMETRY (1 unit = 1 cm), n = 2, 2-norm. Center xbar = (0,0), radius r = 1.2 in
% the first two panels. Third panel: the blob is a closed Catmull-Rom curve
% through (-1.5,0.2), (-0.8,1.3), (0.6,1.4), (1.6,0.5), (1.4,-0.9), (0.2,-1.4),
% (-1.2,-1.0); its nearest boundary point to the origin is at distance > 1.2, so
% the dashed circle of radius delta = 0.6 lies inside it.
%
% GREYSCALE: the open ball has a DASHED boundary, the closed ball a SOLID one,
% both hatched; the difference is also stated in the printed formulas.
\usetikzlibrary{patterns}
\definecolor{okabeblue}{HTML}{0072B2}

\begin{tikzpicture}[
    >={Latex[length=2.0mm]},
    font=\small,
    pt/.style={fill=black, draw=black, circle, inner sep=0pt, minimum size=3.2pt},
    lab/.style={font=\scriptsize, inner sep=1.5pt},
    panelcap/.style={font=\small, align=center, text width=48mm, anchor=north},
  ]
% ------------------------------------------------------------- open ball
\begin{scope}
  \path[pattern=north east lines, pattern color=okabeblue!55] (0,0) circle (1.2);
  \draw[okabeblue, line width=1.2pt, dashed] (0,0) circle (1.2);
  \node[pt] at (0,0) {};
  \node[lab, anchor=north east] at (-0.03,-0.04) {$\bar{\mathbf{x}}$};
  \draw[->, line width=0.7pt] (0,0) -- (0.85,0.85) node[midway, above left, lab] {$r$};
  \node[panelcap] at (0,-1.4) {\textbf{open ball}\\[-1pt]
        $\{\mathbf{x} : \lVert \mathbf{x} - \bar{\mathbf{x}} \rVert < r\}$};
\end{scope}
% ------------------------------------------------------------- closed ball
\begin{scope}[xshift=4.6cm]
  \path[pattern=north east lines, pattern color=okabeblue!55] (0,0) circle (1.2);
  \draw[okabeblue, line width=1.4pt] (0,0) circle (1.2);
  \node[pt] at (0,0) {};
  \node[lab, anchor=north east] at (-0.03,-0.04) {$\bar{\mathbf{x}}$};
  \draw[->, line width=0.7pt] (0,0) -- (0.85,0.85) node[midway, above left, lab] {$r$};
  \node[panelcap] at (0,-1.4) {\textbf{closed ball}\\[-1pt]
        $\{\mathbf{x} : \lVert \mathbf{x} - \bar{\mathbf{x}} \rVert \le r\}$};
\end{scope}
% ------------------------------------------------------------- neighborhood
\begin{scope}[xshift=9.4cm]
  \path[pattern=north east lines, pattern color=black!30]
     [smooth cycle, tension=0.8] plot coordinates
     {(-1.5,0.2) (-0.8,1.3) (0.6,1.4) (1.6,0.5) (1.4,-0.9) (0.2,-1.4) (-1.2,-1.0)};
  \draw[black!70, line width=0.9pt, dashed]
     [smooth cycle, tension=0.8] plot coordinates
     {(-1.5,0.2) (-0.8,1.3) (0.6,1.4) (1.6,0.5) (1.4,-0.9) (0.2,-1.4) (-1.2,-1.0)};
  \draw[okabeblue, line width=1.2pt, dashed, fill=white] (0,0) circle (0.6);
  \node[pt] at (0,0) {};
  \node[lab, anchor=north east] at (-0.03,-0.04) {$\bar{\mathbf{x}}$};
  \draw[->, line width=0.7pt] (0,0) -- (0.42,0.42) node[midway, above left, lab] {$\delta$};
  \node[lab, anchor=south west] at (0.85,0.85) {$\mathcal{N}$};
  \node[panelcap] at (0,-1.55) {\textbf{a neighborhood} $\mathcal{N}$ of $\bar{\mathbf{x}}$\\[-1pt]
        contains an open ball $\mathcal{B}(\bar{\mathbf{x}},\delta)$};
\end{scope}
\end{tikzpicture}
